\section{Conclusion}
\label{sec:conclusion}

This paper studied how a controller should publish a multi-field record
to a supervisory monitor over a coherent fabric. The main result is
simple. Coherence is the substrate, not the publication contract. Both
direct sharing and DMA mirroring fail when the record is published
without an explicit primitive, and both become correct once one is
imposed. Among the correct designs in this study, generation-based
double buffering is the strongest point in the design space: it delivers
zero torn reads, the lowest CPU-side coherence traffic, the shortest
whole-workload simulated time, and negligible reader retries. The
practical lesson is to specify the publication unit, the validity check,
and the reader's behavior on overlap. A coherent fabric can move the
latest line quickly, but it cannot tell the reader whether several
fields belong to one logical publication.

Although the case study uses PWM-inspired witness and evidence streams,
the result applies more broadly to supervisory channels that export a
fixed-size multi-field control or state summary. Our evidence is
protocol-level rather than cycle-accurate, because the platform is gem5
ARM SE mode with Ruby MESI\_Two\_Level, and the timing metric is
whole-workload simulated time rather than a per-publication latency
distribution. Within that scope, the conclusion is stable: monitor
channels should be designed around a publication primitive, not around
coherence alone, and primitive choice can dominate transport choice even
when every design runs over the same coherent substrate.
